% ================================================
% TOGGLE: comment one, uncomment the other
% ================================================
\documentclass[aps, pre, twocolumn, 10pt, groupedaddress, superscriptaddress, linenumbers]{revtex4-2}
% \documentclass[10pt, letterpaper, onecolumn]{article}

% ------------------------------------------------
% BASIC ENCODING AND LANGUAGE SETTINGS
% ------------------------------------------------
\usepackage[utf8]{inputenc}
\usepackage[english]{babel}
\usepackage{csquotes}
\usepackage{xpatch}

% ------------------------------------------------
% PAGE LAYOUT AND SPACING
% ------------------------------------------------
% \usepackage[top=1in, bottom=1in, left=0.75in, right=0.75in]{geometry} % article only — revtex manages margins
\usepackage{setspace}
% \usepackage[switch]{lineno} % article only — incompatible with revtex4-2
% \linenumbers % article only — incompatible with revtex4-2

% ------------------------------------------------
% AUTHOR / TITLE BLOCK
% ------------------------------------------------
% \usepackage{authblk}   % article only — incompatible with revtex4-2

% ------------------------------------------------
% MATHEMATICS
% ------------------------------------------------
\usepackage{amsfonts, amssymb, amsmath}
\usepackage{mathtools}
\usepackage{xfrac}
\usepackage{cancel}

\DeclareMathOperator*{\argmax}{argmax}
\DeclareMathOperator*{\argmin}{argmin}

% ------------------------------------------------
% TABLES AND FLOATS
% ------------------------------------------------
\usepackage{array}
% \usepackage{float} % article only — incompatible with revtex4-2

% ------------------------------------------------
% GRAPHICS AND COLORS
% ------------------------------------------------
\usepackage[dvipsnames]{xcolor}
\usepackage{graphicx}
\graphicspath{{figures/}}
% \usepackage{subcaption}
% \usepackage{caption}
% \captionsetup{
%   aboveskip=5pt,
%   labelfont={bf,footnotesize},
%   labelsep=period,
%   justification=justified,
%   textfont=footnotesize,
%   figurename=Fig.
% }
\makeatletter
\renewcommand{\fnum@figure}{Fig. \thefigure}
\makeatother
% ------------------------------------------------
% CROSS-REFERENCES AND LINKS
% ------------------------------------------------
\usepackage{nameref}
\usepackage{hyperref}
\hypersetup{
  colorlinks=true,
  linkcolor=MidnightBlue,
  citecolor=ForestGreen,
  urlcolor=BrickRed,
  pdfauthor={Your Name},
  pdftitle={Document Title}
}

% ------------------------------------------------
% OPTIONAL UTILITIES
% ------------------------------------------------
\usepackage{changepage}
% \usepackage{lastpage, fancyhdr}  % article only — revtex handles page numbering
% \fancyhf{} % article only — revtex handles page numbering
% \cfoot{\thepage}
% \pagestyle{fancy}
\usepackage{epstopdf}
\usepackage{pifont}
\usepackage[normalem]{ulem}
\usepackage[shortlabels]{enumitem}
\usepackage{tcolorbox}
\newtcolorbox{resultbox}[1][]{
    colback=blue!5!white,
    colframe=blue!50!black,
    fonttitle=\bfseries,
    title=#1
}

% ------------------------------------------------
% BIBLIOGRAPHY
% ------------------------------------------------
% --- revtex4-2: use natbib + apsrev4-2 style ---
\usepackage{natbib}

% ---(place at end of document:)---
  % \bibliographystyle{apsrev4-2}
  % \bibliography{references}

% --- article: use biblatex (comment out natbib block above) ---
% \usepackage[
%   backend=bibtex,
%   style=numeric-comp,
%   sorting=none,
%   maxnames=2,
%   url=false,
% ]{biblatex}
% \DeclareFieldFormat[article]{journaltitle}{#1}
% \addbibresource{references.bib}

% ---(place at end of document:)---
%   \printbibliography

% ------------------------------------------------
% CUSTOM COMMANDS
% ------------------------------------------------
\newcommand{\EQ}{\begin{equation}}
\newcommand{\EE}{\end{equation}}
\newcommand{\EQA}{\begin{eqnarray}}
\newcommand{\EEA}{\end{eqnarray}}
\newcommand{\on}{\text{on}}
\newcommand{\off}{\text{off}}
\newcommand{\mf}{\text{f}}
\newcommand{\act}{\text{act}}
\newcommand{\rec}{\text{rec}}
\newcommand{\step}{\text{step}}
\newcommand{\ens}{\text{ens}}
\newcommand{\peak}{\text{peak}}
\newcommand{\tail}{\text{tail}}
\newcommand{\gen}{\text{gen}}
\newcommand{\largest}{\text{large}}
\newcommand{\nonact}{\text{non-act}}
\newcommand{\p}{p}
\newcommand{\z}{z}
\newcommand{\zbar}{\bar z}
\newcommand{\Z}{{\cal Z}}
\newcommand{\Zbar}{\bar {\cal Z}}
\newcommand{\B}{{\cal B}}
\newcommand{\T}{{\cal T}}
% \newcommand{\DG}{\Delta G}
\newcommand{\DG}{E}
\newcommand{\E}{{\cal E}}
\newcommand{\I}{{\cal I}}
\newcommand{\X}{{\cal X}}
\newcommand{\Po}{{\cal P}}
\newcommand{\K}{{\cal K}}
\renewcommand{\a}{{\bf a}}
\renewcommand{\b}{{\bf b}}
\newcommand{\Ab}{{\bf Ab}}
\renewcommand{\d}{{\rm d}}
\newcommand{\eq}{{\rm eq}}
\newcommand{\nint}{n}
\newcommand{\ep}{\epsilon}
\newcommand{\lam}{\lambda}
\newcommand{\Lam}{\Lambda}
\newcommand{\x}{{\bf r}}
\newcommand{\DKL}{D_{\mathrm{KL}}}

%Colors
\newcommand{\comm}[1]{{\noindent \color{blue} #1}}
\newcommand{\todo}[1]{{\noindent \color{red} #1}}
\newcommand{\out}[1]{{\small \noindent \color{gray} #1}}
\newcommand{\highlight}[1]{{\noindent \color{olive} #1}}
\newcommand{\data}[1]{{\noindent \color{cyan} #1}}%
% \newcommand{\change}[2]{\textcolor{purple}{\sout{#1}} \textcolor{teal}{#2}}
\newcommand{\versions}[2]{\noindent\textcolor{teal}{#1} \noindent\textcolor{purple}{#2}}
\renewcommand{\versions}[2]{\noindent #1}
\newcommand{\new}[1]{\textcolor{teal}{#1}} 

%published script version
%\renewcommand{\comment}[1]{} \renewcommand{\todo}[1]{} \renewcommand{\out}[1]{} \renewcommand{\highlight}[1]{}\renewcommand{\highlight}[1]{} \renewcommand{\data}[1]{} \renewcommand{\change}[1]{}\renewcommand{\new}[1]{}

% Table utilities
\newlength\savedwidth
\newcommand\thickcline[1]{%
    \noalign{\global\savedwidth\arrayrulewidth\global\arrayrulewidth 2pt}%
    \cline{#1}%
    \noalign{\vskip\arrayrulewidth}%
    \noalign{\global\arrayrulewidth\savedwidth}%
}
\newcommand\W[1]{q^{}_{\! #1}}

% Misc
\renewcommand{\epsilon}{\varepsilon}
\parskip 1ex
\newcolumntype{+}{!{\vrule width 2pt}}

% ================================================
% TITLE BLOCK
% ================================================
% --- revtex4-2 style (active) ---
\begin{document}
\title{Exponential Proofreading: A Kinetic Mechanism of B-cell Immunity}
\author{Roberto Morán-Tovar}
\affiliation{Center for Computational Biology, Flatiron Institute, New York, NY, USA}
\affiliation{Institute for Biological Physics, University of Cologne,
 K\"oln, Germany}
\author{Michael Lässig}
\affiliation{Institute for Biological Physics, University of Cologne,
 K\"oln, Germany}

% \date{}

% --- article + authblk style (swap in when using article class) ---
% \author[1]{{\normalsize Roberto Mor\'an-Tovar}}
% \author[1]{{\normalsize Eric Brauns}}
% \author[1]{{\normalsize Michael L\"assig}}
% \affil[1]{{\footnotesize Institute for Biological Physics, University of Cologne,
%   Z\"ulpicherstr.~77, 50937 K\"oln, Germany}}
% \title{Exponential Proofreading: A Kinetic Mechanism of B-cell Immunity}
% \date{}
% \begin{document}


%----------------------------------ABSTRACT--------------------------------------

% --- revtex4-2 ---
\begin{abstract}
Protective B-cell immunity relies on circulating antibodies and on memory B cells that mount rapid recall responses upon reinfection.
Both are seeded by an early expansion of antigen-specific B cells that precedes affinity maturation in germinal centers.
Here, we show that this early phase selects high-affinity clones by a timing mechanism mediated by T-cell help: all activated clones divide at the same rate, but higher-affinity clones start dividing earlier and stop later.
We call this mechanism exponential proofreading.
It predicts a power-law distribution of clone sizes and, because memory B cells pass through it twice, a decay that is twice as steep in recall responses as in primary ones.
Sequencing data from mouse immunization experiments confirm this doubling and are inconsistent with models in which binding affinity sets the division rate.
The second round of proofreading also restricts recall responses to a few high-affinity memory clones, consistent with the observation that most of the memory pool remains unengaged upon boosting.
Finally, the stringency of T-cell help tunes a trade-off between the specificity and the yield of the response, which sets the potency of memory recall.
Together, these results provide a quantitative link between B-cell clonal dynamics and protection against reinfection.
\\
\paragraph*{\bf Significance Statement}
{\small After an infection or vaccination, a few B cells that recognize the pathogen multiply rapidly to produce antibodies and long-lived memory cells.
How the immune system favors the cells that bind best has mostly been studied in germinal centers, where antibodies improve by mutation.
We show that selection already happens earlier, through a different rule: cells that bind better start multiplying sooner and stop later, so they end up as much larger clones.
This rule, which depends on help from T cells, predicts how clone sizes are distributed, why memory responses are dominated by a few large clones, and why most memory cells are not used upon reinfection.
It offers a quantitative basis for predicting protection after infection and vaccination.}
\begin{center}
    \textbf{Keywords}  $|$ B cells $|$ Immune learning  $|$ Proofreading $|$ Exponential dynamics $|$ 
\end{center}
\end{abstract}
\maketitle

% --- article + authblk (comment out block above, uncomment below) ---
% \twocolumn[
% \maketitle
% \begin{abstract}
% ... 
% \paragraph{Keywords}  $|$ B cells $|$ Immune learning  $|$ Proofreading $|$ Exponential dynamics $|$ \\
% \end{abstract}
% ]

\vspace{1cm}


%::::::::::::::::::::::::::::::::::::::::::::::::::::::::::::::::::::::::::::

{\Huge T}he adaptive immune system relies on the coordinated activity of B and T cells to mount effective responses against pathogens.
A fundamental step in this process is the activation and clonal expansion of a small fraction of B cells, which occurs rapidly during the early stages of an immune response~\cite{Lam_2019_Bcells}.
Activated naive B cells, also known as activated precursors (APs), may differentiate into plasmablasts, memory B cells, or germinal center (GC) B cells, each with a distinct role~\cite{pape_different_2011, Lam_2019_Bcells}.
Plasmablasts provide a rapid but short-lived antibody response~\cite{Lam_2019_Bcells}; memory B cells ensure faster and more potent responses upon re-exposure to related pathogens~\cite{Inoue_2024_Memory_review}; and GC B cells undergo rounds of mutation and selection that lead to affinity maturation (AM) of antibodies~\cite{Eisen_2014_AM, Victora_2022_GC}.

In the framework of clonal selection theory~\cite{Burnet_1957_Abs_theory, rajewskyClonalSelectionLearning1996}, the specificity and potency of the B-cell response are determined in two stages: the initial expansion of APs~\cite{pereira_finding_2010} and the subsequent AM in GCs~\cite{Victora_2022_GC}.
As a fast Darwinian process, AM has been the focus of extensive research~\cite{Perelson_1997_Review, AltanBonnet_2020_Review, Victora_2022_GC}, to the point that clonal selection is now often identified with the GC reaction itself.
The early expansion has received far less quantitative attention.

Yet the early expansion matters in its own right.
A large fraction of the memory B-cell pool is generated with little or no AM~\cite{tangye_human_2007, dogan_multiple_2009, Viant_2020_Fate, Mesin_2020_Restricted_clonality, Glaros_2021_Early_memory}.
In particular, a substantial population of early memory B cells arises directly from APs, and its exit from the response is controlled by the decline in antigen availability rather than by AM~\cite{Glaros_2021_Early_memory}.
The properties of this compartment are therefore inherited largely from the initial expansion.
This expansion is itself affinity-dependent and mediated by T-cell help~\cite{Schwickert_2011_Tcells}, which also strongly determines entry into the GC~\cite{yehGerminalCenterEntry2018}.
How this phase selects high-affinity clones, and with what quantitative consequences, is less clear.
In principle, affinity could set how fast clones divide~\cite{Desponds_2021_Population} or for how long; the two alternatives leave distinct signatures in the clonal structure of the responding population, which can be read directly from repertoire sequencing data~\cite{MoranTovar_2024_Primary}.

The early expansion is particularly important in recurrent infections.
In primary infections, the exponential growth of pathogens is capped mainly by innate immunity, independently of the B-cell response~\cite{perelson_hiv-1_1996, baccam_kinetics_2006}.
In a recurrent infection, by contrast, memory B cells must respond rapidly to curb the expanding pathogen~\cite{gray_immunological_1993, palmRemembranceThingsLongTerm2019, Inoue_2024_Memory_review}, by proliferating and differentiating into plasmablasts that produce antibodies in large amounts.
Experiments in mice show that such recall responses are dominated by a few memory clones, while most of the clonal diversity generated by priming is not re-engaged~\cite{Mesin_2020_Restricted_clonality}.
Understanding the early dynamics is therefore essential to explain the specificity, breadth, and timing of recall responses, as well as to characterize the input to GC evolution.

Building on previous work~\cite{MoranTovar_2024_Primary}, we propose that early B-cell expansion selects high-affinity clones by a timing mechanism, which we call \textit{exponential proofreading} (EP): differences in clone size arise not from ``how fast cells divide'' but from ``for how long they divide''.
In our model, T-cell help defines an affinity-dependent window of division: higher-affinity clones acquire sufficient help earlier and lose it later.
EP predicts a power-law distribution of clone sizes and, because memory B cells pass through EP twice, a power law twice as steep in recall responses as in primary ones.
We confirm this doubling in mouse immunization data from two protocols, and show that the measured clone-size statistics are inconsistent with models in which binding affinity sets the division rate.
The second round of EP also restricts which memory clones are re-engaged, consistent with the recall bottleneck observed in prime-boost experiments.
Finally, the stringency of T-cell help tunes a trade-off between the specificity and the yield of the response, which sets the potency of memory recall.

The paper is organized as follows.
We first introduce a model of antigen presentation and T-cell-mediated clonal expansion.
We then derive EP for primary and recall responses, test its predictions against repertoire data, and analyze how T-cell help sets the specificity and potency of the response.
Finally, we discuss the implications of EP for B-cell immunity.


%::::::::::::::::::::::::::::::::::::::::::::::::::::::::::::::::::::::::::::
\section*{Model}
%::::::::::::::::::::::::::::::::::::::::::::::::::::::::::::::::::::::::::::


% -----------------------------------
\begin{figure*}[t]
\begin{center}
\includegraphics[width=.95\linewidth]{model.pdf}
\caption{\textbf{B-cell antigen presentation and T-cell-mediated proliferation model.}
{\bf A} Antigen presentation. Free antigen (orange), proliferating at rate $\lam_A$, binds the BCR of a B cell (purple) with association rate $k_\on$ and dissociation rate $k_\off$; a bound antigen is irreversibly internalized and processed at rate $k_s$, yielding surface pMHC.
{\bf B} B-cell proliferation. The B cell presents pMHC to a cognate T cell (pink) and acquires T-cell help at a net rate $h(\pi)$, set by the surface pMHC $\pi$; the licensed cell then divides at the intrinsic rate $b_B$, giving two daughter B cells.
}
\label{fig:model}
\end{center}
\end{figure*}
% -----------------------------------

\subsection*{Antigen binding and presentation}
B cells are defined by the B-cell receptor (BCR) they express. For a given antigen, BCRs can be characterized by the antigen-BCR binding affinity.
This affinity can be specified by the dissociation constant $K = k_\off/k_\on$, or equivalently by an effective binding energy $\DG = \log\!\left({K}/{K_s}\right)$, with $K_s \equiv {k_s}/{k_\on}$ and where $k_s$ is the effective rate of irreversible internalization, processing and presentation of antigen (see Fig.~\ref{fig:model}A).
Processed antigen is presented on the cell surface as peptide-MHC complexes (pMHCs).
The reference scale $K_s$ is then the dissociation constant at which off-rate and internalization balance, so that $\DG = 0$ marks the crossover between encounters that mostly dissociate ($\DG > 0$) and encounters that mostly lead to productive antigen internalization and presentation ($\DG < 0$).
Larger $\DG$ thus corresponds to weaker binding/lower affinity.

The probability that a single antigen-BCR encounter results in productive internalization is set by the competition between dissociation and processing, $\psi(\DG) = \left(1+\exp{\DG}\right)^{-1},$ which interpolates between $\psi \to 1$ for strong binders ($\DG \ll 0$) and $\psi \to 0$ for weak binders ($\DG \gg 0$).
The rate of pMHC production on the surface of a B cell with binding energy $\DG$ is therefore proportional to the antigen concentration, the per-cell association rate, and the internalization probability, $u_p(t, \DG) = k_\on\, N_A(t)\, \psi(\DG)$, where $k_\on$ is the effective total association rate between antigen and the BCRs of a single B cell.

\subsection*{Two-stage proliferation}
B-cell division proceeds through two sequential stages.
First, the B cell must accumulate sufficient pMHC and engage cognate T cells to acquire a proliferation signal~\cite{Schwickert_2011_Tcells}; this occurs at an effective rate $h(t, \DG)$ (Fig.~\ref{fig:model}B).
Second, the B cell undergoes division at an intrinsic rate $b_B$ set by the duration of a cell cycle (Fig.~\ref{fig:model}B).
% \EQ
% B_i^o + T^o \xrightarrow{\,h(t,\DG)\,} B_i^* + T^* \xrightarrow{\,b_B\,} 2\,(B_i^o + T^o).
% \nonumber
% \EE
Because the two stages are sequential, the mean time per cell division is the sum of the two waiting times, giving an effective division rate of harmonic-mean form
\EQ
\lam_B(t, \DG) = \left(\frac{1}{h(t, \DG)} + \frac{1}{b_B}\right)^{-1}.
\label{eq:lambdaB}
\EE
The slower of the two stages always dominates: $\lam_B \approx h$ when T-cell help acquisition is rate-limiting, and $\lam_B \approx b_B$ when the intrinsic cell division is rate-limiting.

The net help rate $h$ depends on the presented surface pMHC density and on the availability of free T cells.
We take it proportional to the pMHC density, $h(t, \DG) = h_0\,\pi(t, \DG)$, where $h_0$ is the rate at which a B cell acquires help per presented pMHC complex.
It lumps together the size of the cognate T-cell pool, the co-stimulatory requirements of the T-B interaction, and the efficiency of pMHC-TCR engagement, and is a property of the host rather than of the responding clone.

Two features of the response make T-cell availability a non-limiting factor in our regime.
First, cognate T cells proliferate in parallel with B cells~\cite{garsideVisualizationSpecificLymphocyte1998}, replenishing the engaged pool on the same timescale at which it is depleted.
Second, because a single antigen presents many peptide epitopes to T cells but only a few B-cell–accessible surface epitopes, the pool of cognate T cells available to help is broad relative to the small number of B-cell clones initially competing for that help at the onset of the response~\cite{besavillaPreexistingCD4Cell2023}.

\subsection*{T-cell mediated B-cell clonal expansion}
We describe the dynamics of the immune response through three coupled variables: antigen concentration $N_A(t)$, pMHC surface density per B cell $\pi(t, \DG)$, and B cell clone size $N_B(t, \DG)$ (See Materials and Methods).

The two-stage proliferation scheme above is intrinsically stochastic at the single-cell level, with affinity-dependent waiting times for antigen acquisition and division.
Here we focus on the mean behavior at the population-level, treating clone size $N_B(t, \DG)$ as a continuous variable that grows deterministically and the density of pMHC per-cell $\pi(t, \DG)$ as a smooth function of time.

The harmonic-mean structure of $\lambda_B$ ensures that the slowest process dominates the division rate.
This sets a characteristic density of pMHC at which T-cell help becomes likely on the intrinsic division time scale of B-cells, defined by $h(\pi_c) = b_B$, or equivalently $\pi_c = b_B/h_0$.
Recent experimental evidence has shown that an activation threshold for T cells is directly determined by the cooperative activity of pMHC-TCR to trigger downstream signaling~\cite{wither_antigen_2023, white_proofreading_2025}.
Additionally, the threshold can be an indirect consequence of other mechanisms such as quorum detection~\cite{Butler_2013_Tcells_QS} or effective aggregate dwell times~\cite{govern_fast_2010}.
In our model, this threshold emerges from the mean-field dynamics, giving rise to three phases of clonal expansion (Fig.~\ref{fig:ep}).

\textit{Phase 0: Antigen-limited.} At low antigen concentration, the per-cell pMHC density is below $\pi_c$ and T-cell help is unlikely on the timescale of a cell division, $h \ll b_B$. Clonal expansion is negligible.

\textit{Phase 1: Division-limited with antigen growth.} As $N_A(t)$ grows, $\pi$ rises above $\pi_c$ and the bottleneck shifts to the intrinsic cell division, $\lambda_B \approx b_B$. All activated clones expand at the same rate, independently of affinity.

\textit{Phase 2: Division-limited with antigen decay.} After the antigen peak at $t_p$, $\pi$ declines, and clonal expansion continues at $\lambda_B \approx b_B$ until $\pi$ falls below $\pi_c$.

The transitions between phases occur at affinity-dependent times, which we derive in the Results.

\subsection*{Density of B cell clones}
The B-cell repertoire consists of a large number of different clones.
We define the density $\Omega_0(\DG)$, where $\Omega_0(\DG)\,\d\DG$ is the expected number of naive clones with binding energy in $[\DG,\DG+\d\DG]$. 
At the tail of high-affinity clones available in the naive repertoire, we assume $\Omega_0(\DG) =\omega_0 e^{\beta^*\DG}$, reflecting the exponential increase in the number of clones of lower affinity.
The exponent $\beta^*$ is a summary measure of repertoire complexity~\cite{MoranTovar_2024_Primary}.
We assume that naive repertoires are randomly sampled from an underlying amino acid sequence distribution~\cite{Elhanati_2015_Diversity} so that most lineages bind only weakly to any given antigen ($\DG \gg 1$).
Following previous work, we assume that the internalization rate is tuned to the best binder of the naive repertoire, $k_s/k_{\rm on} \approx K^*$ (equivalently $\DG^* \approx 0$), so that activation efficiency saturates at the high-affinity edge of the naive repertoire~\cite{MoranTovar_2024_Primary}.
We note that $\DG^*\approx 0$ is substantially above the global minimum $\DG_m$ (the Master sequence of the binding motif).

% -----------------------------------
\begin{figure}[t]
\begin{center}
\includegraphics[width=\linewidth]{ep.pdf}
\caption{\textbf{Exponential proofreading: clonal expansion as a timing problem.}
Dynamics of a primary (left) and recall (right) response for three clones of decreasing binding affinity (dark to light green).
The phases of clonal expansion are shown with grey shadowing.
{\bf A} Antigen load $N_A(t)$ grows exponentially at rate $\lam_A$ until the peak at $t_p$, then decays at rate $\delta_A$.
{\bf B} Per-cell presented-antigen density $\pi(t,\DG)$ for a clone with $\DG^*$ (darker green line) and two other lower affinity clones ($\DG =$ \todo{$x$ and $x$}).
{\bf C} Clone size $N_B(t,\DG)$ for three clones (same color code and affinities as in {\bf B}).
Each clone expands at full speed only while $\pi > \pi_c$, bracketed by the activation front (open diamond, $t_\on$) and the termination front (filled diamond, $t_\off$).
Higher-affinity clones cross $\pi_c$ earlier and fall below it later, dividing over a longer interval $\Delta t(\DG)$ and reaching a larger final size.
{\bf D} Moving fronts of activation (inverse of $t_{\on}(\DG)$; solid lines) and termination (inverse of $t_{\off}(\DG)$; dashed lines) for the full model (green) and the null model without T-cell help (grey).
}
\label{fig:ep}
\end{center}
\end{figure}
% -----------------------------------

\subsection*{Diversity and specificity of responses}
The clonal population of B cells can be studied at two distinct levels.
Using a (usually) genetic definition of clone, we can describe the total B-cell population with the discrete distribution $\{p_i\}$ where $p_i$ is the fraction of B cells in the clone $i$.
For the naive repertoire, we assume that each clone has initially clone size 1 and therefore $p_i^0 = 1/L_0$, where $L_0$ is the total number of different naive B cells.
After an immune response, the clone size distribution is modified due to clonal expansion, leading to a new clone-size distribution $p_i\neq p_i^0$.

Alternatively, we can group B cells by their binding affinity for a fixed antigen and define the pdf $q(\DG)$ such that the fraction of B cells with energy in the interval $[\DG, \DG + \d\DG]$ is given by $q(\DG)\d\DG$.
First, we define the naive pdf $q_0(\DG) = \Omega_0(\DG)/L_0$ where $\Omega_0(\DG)$ is the density of the naive repertoire (see grey curve in Figs.~\ref{fig:potency}A and B).
Similarly, we define a pdf for the binding energies of the responding repertoires after clonal expansion, $q(\DG)$ (see colored curves in Figs.~\ref{fig:potency}A and B).
% For the analytical results, we assume that binding energies of the naive and the responding repertoires are Gaussian-distributed with universal variance $\sigma^2$ and a sharp high-affinity bound at $\DG^*$, and with a repertoire-dependent mean $\mu_0$.
% We characterize the naive repertoire by the affinity scale $\Delta\DG \equiv \mu_0 - \DG^*$, the gap between the mean binding energy and that of the best clone.
% This scale does not depend on the choice of energy origin; with our convention $\DG^* = 0$, it equals the mean, $\Delta\DG = \mu_0$.

% In the mean-field limit, we assume that for responding repertoires there is a sharp low-affinity bound at $\DG_{\max} = \DG_\on(t_p)$, the largest reach of the moving front of activation during the response.
% Altogether, this leads to a double-truncated Gaussian distribution $q(\DG)$ with the same variance $\sigma^2$ but shifted mean $\mu<\mu_0$ (see SI).

We define the diversity of a B-cell immune response as the Shannon entropy of the discrete clone-size distribution $S_c[\{p_i\}]$.
Additionally, we define the specificity of a B-cell immune response as the relative entropy between the responding and the naive distributions, $\DKL[q\|q_0]$, and the corresponding effective fraction of the naive repertoire engaged by the response as $\phi \equiv \exp\left[-\DKL[q\|q_0]\right]$.
As will become clear in the next sections, the clonal diversity $S_c$ depends on the definition of a clone, while the specificity does not.

\subsection*{Rapidity and potency of responses}
To study the real-time binding capacity of the ensemble of secreted antibodies, we define a time-dependent partition function $Z(t) = \sum_i c_oN_B(t,\DG_i) \exp[-\DG_i]$, where the sum runs over all activated and proliferating B-cell clones and with $c_o$ a constant that accounts for differentiation efficiency, per-cell antibody secretion rate (See SI).
Each time a new clone with affinity $\DG$ is activated, a term $\exp[-\DG]$ is added to the sum.
Additionally, as activated clones proliferate, their contribution to the sum increases.
Therefore, the partition function increases during an immune response due both to activation and to clonal proliferation.
From $Z(t)$, we can calculate the probability $p_b(Z)$ that a viral particle is bound by antibodies assuming thermodynamic equilibrium ~\cite{phillips_physical_2019} (See SI).
In a minimal model for viral fitness where its growth rate depends on $p_b$ (see Materials and Methods), we define a critical value $Z_c$ at which pathogen growth is neutralized ($p_b\sim \mathcal O(1)$).

We define the control time $t_c$ as the first passage of $Z$ to $Z_c$, $Z(t_c) = Z_c$.
Since the antigen grows at rate $\lam_A$ until control, we define the potency of the real-time recall response as the log-fold reduction in viral-load peak,
\EQ
P \equiv \lam_A\,(t_p^o - t_c),
\label{eq:P}
\EE
where $t_p^o$ is the viral-load peak of a naive infection.


%::::::::::::::::::::::::::::::::::::::::::::::::::::::::::::::::::::::::::::
\section*{Results}
%::::::::::::::::::::::::::::::::::::::::::::::::::::::::::::::::::::::::::::

\subsection*{Exponential Proofreading: Clonal expansion as a timing problem}
A clone with binding energy $\DG$ is activated when its pMHC density reaches the threshold, $\pi(t,\DG)=\pi_c$, marking the transition from Phase 0 to Phase 1.

\textit{Affinity-graded activation.}
Before activation, $\pi$ tracks quasi-statically the antigen load, $\pi \propto N_A(t)\,\psi(\DG)$ (see Materials and Methods).
The activation condition therefore defines a moving front $\DG_\on(t)$ in energy space: by time $t$, all clones with $\DG < \DG_\on(t)$ have been activated.
For exponential antigen growth and $\psi(\DG)\approx \exp{[-\DG]}$, the front advances linearly in time with velocity $v_\on=\lam_A$, because advancing it by $\delta\DG$ requires $N_A$ to grow by a factor $\exp{[\delta\DG]}$, which always takes the same time $\delta\DG/\lam_A$.
The activation time of a clone with binding energy $\DG$ is
\EQ
t_\on(\DG) = t^* + \frac{\DG}{v_\on},
\label{eq:t_on}
\EE
where $t^* = t_\on(\DG^*)$ is the activation time of the host's best clone (Fig.~\ref{fig:ep}D, solid green lines).
The front velocity is thus set by the pathogen growth rate.
In previous work, the front velocity was also shaped by a molecular kinetic proofreading scheme, which set the specificity of the response~\cite{MoranTovar_2024_Primary}.
Here, as we show below, specificity is set by T-cell help.
At the antigen peak $t_p$, the front stalls at $\DG_{\max} \equiv \DG_\on(t_p)$, and the total number of activated lineages is $L_\act \sim \exp{[\beta^* \DG_{\max}]}$.

\textit{Affinity-graded termination.}
After the peak, $\pi$ decays at rate $v_\text{off} = \min(\delta_A, \lambda_B)$, set by the slower of two channels: dilution by cell division and decay of the antigen source, assuming slow surface turnover $\delta_\pi \ll \lambda_B$ (see Materials and Methods and SI).
Because high-affinity clones reach a higher pMHC level at the peak, the condition $\pi = \pi_c$ defines a second front $\DG_\off(t)$ that sweeps in the opposite direction: weak binders exit first, strong binders last.
The termination time of a clone with binding energy $\DG$ is 
\EQ
t_\off(\DG) = t_p + \frac{\DG_{\max} - \DG}{v_\off}
\label{eq:t_off}
\EE
where now the affinity-dependence is negative (Fig.~\ref{fig:ep}D, dashed green lines).

\textit{Affinity-dependent clonal expansion.}
Between $t_\on$ and $t_\off$, all activated clones divide at the same rate $\lam_B \approx b_B$, so the final clone size is $N_B(\DG) = \exp[\lam_B \Delta t(\DG)]$, with $\Delta t(\DG) = t_\off(\DG) - t_\on(\DG)$ (Fig.~\ref{fig:ep}C).
Because both fronts are linear in $\DG$, so is $\Delta t(\DG)$ (see Materials and Methods), which gives the central EP relation between final clone size and binding energy,
\EQ
N_B(\DG) = e^{-\alpha(\DG-\DG_{\max})},
\quad
\alpha = \lam_B\left(\frac{1}{v_\on} + \frac{1}{v_\off}\right).
\label{eq:alpha}
\EE

Three features of this result deserve emphasis.
First, all activated clones proliferate at the same rate, regardless of binding energy; what differs between clones is not how fast they divide, but for how long.
Second, Eq.~\ref{eq:alpha} is essentially kinematic: it requires only that both fronts are linear in $\DG$ and that clonal expansion is exponential.
Third, the threshold $\pi_c$ enters $t_\on$ and $t_\off$ identically, through $t^*$ and $\DG_{\max}$, and drops out of the exponent $\alpha$.
The two-front picture thus converts the affinity spectrum of the naive repertoire into a hierarchy of clone sizes, whose shape is set by three macroscopic rates: the B-cell division rate and the two front velocities.


\begin{figure*}[t]
\begin{center}
\includegraphics[width=\linewidth]{csd.pdf}
\caption{\textbf{Clone-size statistics of primary and recall responses.}
{\bf A} Clone-size profiles for three representative mice each from primary (orange) and recall (blue) responses; bars show clone size for clones sorted by clone-size rank, on a log scale. Recall profiles are visibly less uniform, reflecting the steeper clone-size hierarchy predicted by two rounds of EP.
{\bf B} Clone-size-rank curves, normalized to the largest clone $N_1$ and pooled across mice. Faint lines are individual mice; symbols are cohort averages for primary (orange stars) and recall (blue squares). Solid lines are Zipf fits $N_k/N_1 \sim k^{-\zeta}$, giving $\zeta = 0.51 \pm 0.04$ (primary) and $\hat\zeta = 1.13 \pm 0.03$ (recall), a ratio $\hat\zeta/\zeta = 2.2 \pm 0.2$ consistent with the predicted exponent doubling.
{\bf C} Mouse-resolved scaling between the number of activated lineages $L_\act$ and the largest clone size $N_1$. Each point is one mouse, coded by phenotype (GC, orange; GC$+$fm, blue) and experiment (markers). Dashed curves are the predicted relations $L_\act \sim N_1^{1/\zeta}$ using the cohort-averaged $\zeta$ from {\bf B}, with no additional fit parameters; the black dashed line marks the $\zeta = 1$ separatrix dividing the diversity-dominated ($\zeta < 1$, primary) and size-dominated ($\zeta > 1$, recall) regimes. Shading indicates the two regimes.
}
\label{fig:csd}
\end{center}
\end{figure*}
% -----------------------------------

\subsection*{Memory response as a second round of EP}
Memory B cells produced during a primary response provide a head start for subsequent immune challenges in two ways: they are present at an elevated abundance compared to naive clones~\cite{Mesin_2020_Restricted_clonality}, and they have a lower activation threshold than naive cells~\cite{tangyeMemoryCellsEffectors2009, wongBasicsMemoryBcell2019, Inoue_2024_Memory_review}.
Based on that, here we propose a model of memory production after a primary response and its reactivation during a recall response.

We assume that at the end of the primary response, each B cell commits to the memory fate with an affinity-independent probability $\gamma$ and the memory pool has a global carrying capacity $M$, which admits a fraction $\nu = M/(\gamma Y + M)$ of the committed cells, where $Y = \sum_i N_B(\DG_i)$ is the primary yield.
The total number of memory cells is then $\hat Y_o = \nu\gamma Y$ (we use $\hat\cdot$ to denote quantities in the memory compartment).
In the saturated regime ($\gamma Y \gg M$), $\nu \approx M/(\gamma Y)$ and the memory pool saturates at $\hat Y_o \approx M$; in the unsaturated regime ($\gamma Y \ll M$), $\nu \approx 1$ and memory tracks the primary yield, $\hat Y_o \approx \gamma Y$.
In both regimes, memory formation is a uniform rescaling of the primary response: every cell enters memory with the same probability $\nu\gamma$, regardless of its affinity.
Large clones therefore keep their relative sizes, and memory inherits the affinity grading of the primary expansion, $\hat N_B^o(\DG) = \nu\gamma N_B^*\, e^{-\alpha\DG}$.
This holds as long as the best clones are retained, $\nu\gamma N_B^* \gtrsim 1$; otherwise, sampling also thins the high-affinity edge (see Materials and Methods).

Upon re-challenge, memory clones undergo the same kinetic readout described above.
Their lower activation threshold shifts the timing of activation, but not the exponent (as shown above, $\pi_c$ drops out of $\alpha$).
The activation and termination fronts therefore advance with the same velocities ($v_\on$ and $-v_\off$, respectively) as in the primary response, and the proliferation time $\Delta\hat t(\DG)$ has the same affinity dependence (see Materials and Methods).

Thus, the final size of a memory clone with binding energy $\DG$ is given by
\EQ
\hat N_B(\DG) \sim e^{-\hat\alpha\DG},
\qquad
\hat\alpha = 2\alpha .
\label{eq:alpha2}
\EE

The memory clone-size distribution is steeper than the primary one because two rounds of EP compound. 
The new affinity-graded proliferation during memory reactivation occurs on top of the grading inherited from the primary expansion.
This doubling requires no change in the molecular machinery of the recall response.

The second round of EP also restricts which memory clones take part.
In a recall challenge with the same initial antigen load and growth rate as the primary one, the activation front stalls at the control time $t_c$, when the antigen load is lower than the primary peak by a factor $e^{-P}$ (Eq.~\ref{eq:P}).
Its reach is therefore $\hat\DG_{\max} = \DG_{\max} - P + \Delta$, where $\Delta \equiv \log(\pi_c/\hat\pi_c)$ is the activation advantage of memory cells (see Materials and Methods).
Because the memory pool, like the naive repertoire, is dominated by low-affinity lineages, only a fraction $\hat L_\act/L_\act \approx e^{-\beta^*(P - \Delta)}$ of them is re-engaged (Eq.~\ref{eq:bottleneck}).
A protective recall response is thus carried by a few high-affinity memory clones while most of the memory pool remains unengaged, consistent with the clonality bottleneck observed in prime-boost experiments~\cite{Mesin_2020_Restricted_clonality}.

\subsection*{EP explains B-cell clonal structure upon recurrent immunizations}
The EP mechanism predicts that clone size decays exponentially with binding energy as given by Eq.~\ref{eq:alpha}.
Combined with the repertoire density $\Omega_0(\DG) \sim \exp[{\beta^*\DG}]$, this predicts the statistical structure of the responding population, an observable that can be extracted directly from repertoire sequencing data~\cite{MoranTovar_2024_Primary}.
In the mean-field deterministic limit, the clone-size-rank relation, obtained by ordering clones from largest $(k=1)$ to smallest $(k = L_\act)$, is a Zipf's law, $N_k \sim k^{-\zeta}$~\cite{MoranTovar_2024_Primary}.

Within a single response, the power-law form is a signature of the timing mechanism.
In the naive repertoire, lower-affinity clones are exponentially more abundant; EP makes them exponentially smaller in size upon clonal expansion, and the two competing exponential functions in binding energy yield a power law.
Below, we contrast this with models in which affinity sets the expansion rate.

Using Eqs.~\ref{eq:alpha} and~\ref{eq:alpha2}, the clone-size exponents for primary and memory responses are
\EQA
\zeta &=& \frac{\alpha}{\beta^*} \qquad \text{primary}\nonumber \\
\hat\zeta &=& 2\zeta \,\qquad \text{recall}.
\label{eq:zetas}
\EEA
The compounding of two identical rounds of EP therefore predicts the ratio $\hat\zeta/\zeta = 2$, independent of all model parameters.
We note that these predictions concern the early, pre-GC phase of one primary and one recall response, and assume that the memory pool inherits its clonal structure from the primary expansion with little or no affinity maturation~\cite{tangye_human_2007, dogan_multiple_2009, Viant_2020_Fate, Mesin_2020_Restricted_clonality, Glaros_2021_Early_memory}.

To test this result of the model, we analyzed published data from mouse immunization experiments~\cite{Mesin_2020_Restricted_clonality}.
In these experiments, naive mice were immunized one or two times with the same antigen (Materials and Methods), generating in this way primary and recall immune responses.
A few days post-immunization (p.i.), single-cell sequencing analysis was performed, obtaining in this way a fine-grained characterization of the clonal composition of the responding B-cell population.
Among the distinct B-cell sub-populations that were sequenced, we focused our analysis on early GC B cells (day 6 p.i.), when GCs have just formed from the initial clonal expansion, and when GC B cells carry few or no somatic hypermutations.
Importantly, in all primary responses, responding B cells were fate-mapped.
In this way, during recall responses, it was possible to identify if a responding B cell belonged to the memory pool and had participated in the primary response (fate-mapped) or was a naive B cell before \textit{de novo} activation.

In Fig.~\ref{fig:csd}A we show some examples of the clonal composition of activated B-cell populations.
We show the 10 largest clones in each immunized mouse (the complete data is shown in SI).
Consistent with our model, we see a different degree of concentration of B cells between primary and recall responses.
Pooling three experiments sharing the same prime/boost protocol with antigen CGG-alum, we measure $\zeta = 0.51 \pm 0.04$ for primary responses and $\hat\zeta = 1.14 \pm 0.06$ for recall responses, giving the ratio $\hat\zeta/\zeta = 2.2 \pm 0.2$, consistent with the model prediction.
A fourth experiment using a distinct protocol (infection with Influenza followed by boost with recombinant HA) gives $\zeta = 0.50 \pm 0.09$ and $\hat\zeta = 1.15 \pm 0.04$, with the same ratio within errors (Fig.~\ref{fig:csd}B).
The doubling prediction therefore holds across two distinct immunization protocols.
Cohort-level $\zeta$ values are reproducible across experiments; individual-mouse fits show larger spread, especially for memory responses (see transparent lines in Fig.~\ref{fig:csd}B).

The Zipf fit characterizes the average clone-size hierarchy within a cohort but collapses mouse-to-mouse variation into a single curve.
A complementary, mouse-resolved test of the EP picture comes from the scaling relations between observables that vary between mice within a cohort.
Three quantities can be extracted directly from each mouse's repertoire data: the largest clone size $\tilde N_1$, the total cell count $\tilde Y$, and the number of activated lineages $\tilde L_\act$ (we use $\tilde \cdot$ to denote quantities in a given sub-sample of the total expanded population).
Under the assumption that $\alpha$ and $\beta^*$ are shared among mice, while $\DG_{\max}$ varies between mice (e.g.\ due to dose, innate protection, or individual repertoire variability), Zipf's law predicts power-law relations between these three quantities, with exponents fixed by the same $\zeta$.

The largest clone and the number of activated lineages are set by the reach of the activation front.
The largest clone is the best binder, $N_1 = N_B(\DG^*) \sim \exp[\alpha\DG_{\max}]$ (Eq.~\ref{eq:alpha} with $\DG^* \approx 0$), and the number of activated lineages is $L_\act \sim \exp[\beta^*\DG_{\max}]$.
Eliminating $\DG_{\max}$ gives
\EQ
L_\act \sim N_1^{1/\zeta}.
\label{eq:scaling}
\EE

The same relation holds for the subsampled quantities (see SI).
The relation in Eq.~\ref{eq:scaling} is convex for $\zeta < 1$ and concave for $\zeta > 1$.
The value $\zeta = 1$ also separates two regimes of the total yield, $Y = \sum_k N_k$: for $\zeta < 1$, $Y$ is dominated by the many small clones near $\DG_{\max}$ (diversity-dominated); for $\zeta > 1$, by the few largest clones near $\DG^*$ (size-dominated).
The two response types in our data sit on opposite sides of this transition: primary responses in the diversity-dominated regime ($\zeta < 1$, see orange in Fig.~\ref{fig:csd}C), and memory responses in the size-dominated regime ($\zeta > 1$, see blue in Fig.~\ref{fig:csd}C).
The crossover marks the point where the clone-size decay exponent $\alpha$ exactly balances the repertoire exponent $\beta^*$.

The mouse data span more than an order of magnitude in $N_1$ and $L_\act$, and follow the predicted scaling behavior using the cohort-averaged $\zeta$ values from the Zipf analysis (see dashed lines in Fig.~\ref{fig:csd}C).
The scaling test is a non-trivial cross-check of the EP picture: the exponents observed in the mouse-resolved scaling clouds are predicted by the same $\zeta$ extracted from cohort-averaged Zipf fits, with no additional fit parameters.

The measured exponents can be cross-checked against independent estimates of the rates $\lam_A$ and $\lam_B$.
Typical B-cell doubling times during primary responses are $\sim 12-24$~hours, corresponding to $\lam_B \sim 1.0-2.0$~day$^{-1}$~\cite{pape_different_2011, mesin_germinal_2016, Victora_2022_GC, gitlin_clonal_2014}; typical acute viral growth rates span $\lam_A \sim 3-6$~day$^{-1}$~\cite{perelson_hiv-1_1996, baccam_kinetics_2006}, and effective antigen growth rates after vaccination are comparable~\cite{MoranTovar_2024_Primary}.
The ratio $\lam_B/\lam_A$ therefore lies in the range $0.17-0.67$, and in the dilution-limited regime, $v_\off = \lam_B$, Eq.~\ref{eq:alpha} gives $\alpha = 1 + \lam_B/\lam_A \approx 1.2-1.7$.

Inverting Eq.~\ref{eq:zetas}, the measured exponents constrain the repertoire complexity: the primary relation gives $\beta^* = \alpha/\zeta \sim 2.3-3.3$, and the recall relation gives $\beta^* = 2\alpha/\hat\zeta \sim 2.1-2.9$.
Both estimates are consistent with numerical estimates of the repertoire complexity derived in previous work for different binding energy models~\cite{MoranTovar_2024_Primary}.
Moreover, in a Gaussian approximation for the density of naive B cells, assuming BCRs and antigens are sequences of length $\ell$ amino acid residues, and that the total binding energy is a linear sum over the $\ell$ BCR–antigen residue contacts that form the paratope, the repertoire exponent is related to the naive repertoire size and binding motif length as $\beta^* = \sqrt{2\log{L_0}/\ell}$ (see Materials and Methods).
Setting $L_0$ to the estimated total number of naive B cells in mice $L_0 \approx 10^8$~\cite{AltanBonnet_2020_Review}, the inferred $\beta^*$ is consistent with paratopes of length $\ell = 4-9$, in line with estimates of the number of antibody residues that make direct contact with antigens~\cite{akbarCompactVocabularyParatopeepitope2021}.

% -----------------------------------
\begin{figure*}[t]
\begin{center}
\includegraphics[width=\linewidth]{potency.pdf}
\caption{\textbf{Potency and specificity of B-cell responses across the T-cell help threshold.}
{\bf A} Cell binding-energy densities after a primary response for the total expanded B-cell population (dashed) and the sub-sampled memory formed from it (solid). Grey denotes the naive density $\Omega_0(\DG)$; we show three pair of curves for $\pi_c = 3.2, 32.$ and $320.$ (red, green and cyan, respectively)
{\bf B} Cell binding-energy densities after a recall response for the total expanded B-cell population (colors as in {\bf A}).
{\bf C} Partition function $\hat Z(t)$ for a set of recall responses for different values of $\pi_c$. Vertical dashed lines denote time of antigen peak under primary response (grey) and under the optimal $\pi_c^*$ threshold value in a recall response (black). Filled markers highlight denote potency $P$ as defined in the main text (squares for recall and circle for a reference primary response). \todo{(discuss properly the naive $Z$ curve)}
{\bf D} Response potency $P$ versus specificity $\DKL[q\|q_0]$, for the primary, recall, and null responses (colors and symbols as in {\bf A} and {\bf B}). 
Other parameters as in Fig~\ref{fig:ep}}
\label{fig:potency}
\end{center}
\end{figure*}
% -----------------------------------

\todo{\subsection*{Differential growth rates cannot reproduce the observed clonal structure}
EP requires that activated clones divide at the saturated rate $b_B$.
The opposite limit of Eq.~\ref{eq:lambdaB}, in which T-cell help remains rate-limiting throughout the response ($\pi \ll \pi_c$, so that $\lam_B \approx h_0\pi$), is the natural alternative: affinity then sets how fast clones divide, not for how long.
Because $\pi \propto \psi(\DG)$ at all times in this limit, integrating Eq.~\ref{eq:NB} gives
\EQ
\log N_B(\DG) \approx \mathcal{A}\,\psi(\DG) \approx \mathcal{A}\, e^{-\DG},
\label{eq:rate_model}
\EE
where $\mathcal{A}$ is independent of $\DG$ and proportional to the integrated antigen exposure (see SI).
This is the simplest model in which affinity sets the expansion rate~\cite{Desponds_2021_Population}.
In EP, $\log N_B$ is linear in $\DG$; here it is exponential in $\DG$.

This difference determines the shape of the rank statistics.
Combining Eq.~\ref{eq:rate_model} with the lineage density $\Omega_0(\DG) \propto \exp{[\beta^*\DG]}$ gives $\log N_k = \log N_1\, k^{-c}$ with $c = 1/\beta^*$ (see SI).
The local Zipf exponent, $\zeta_{\rm eff}(k) \equiv -\d\log N_k/\d\log k = c \log N_k$, decreases with rank, by a factor $10^{-c} \approx 0.37$ per decade of rank for $\beta^* \approx 2.3$, so that rank curves flatten on logarithmic axes.
EP corresponds to $c = 0$ (straight rank curves).
Fitting the rank curves of all four cohorts with a form that interpolates between both cases (SI) gives $c$ consistent with zero ($c < 0.22$ at 95\% bootstrap confidence), and the curvature parameter does not improve the fit ($\Delta{\rm AIC} \approx +2$, the penalty for the additional parameter).
This holds for primary and recall cohorts alike.
For the rate model to account for these straight rank curves, the repertoire exponent would have to exceed $\beta^* = 1/c \approx 4.5$, well above independent estimates of repertoire complexity~\cite{MoranTovar_2024_Primary}.


The two mechanisms also differ in how a second exposure compounds.
In the rate model, memory inherits $\log \hat N_B^o = \log(\nu\gamma) + \mathcal{A}\, e^{-\DG}$ and the recall response adds $\hat{\mathcal{A}}\, e^{-\DG}$, so that log-size differences between clones are amplified by a factor $1 + \hat{\mathcal{A}}/\mathcal{A}$ that depends on the ratio of recall to primary antigen exposure.
A protective recall response reduces this exposure (Eq.~\ref{eq:P}), so that the factor falls below 2.
In EP, the exponent $\alpha$ is set by the front velocities rather than by antigen exposure, and $\hat\zeta/\zeta = 2$ holds independently of the reach of the fronts, the threshold, and the recall antigen load.

Both mechanisms are limits of the same two-stage division rate (Eq.~\ref{eq:lambdaB}); which one operates depends on whether presented antigen exceeds $\pi_c$.
The straight rank curves, together with the exponent ratios $\hat\zeta/\zeta = 2.2 \pm 0.2$ (CGG-alum) and $2.4 \pm 0.5$ (influenza/HA) (Fig.~\ref{fig:csd}B), place early B-cell expansion in the division-limited, timing regime rather than in a regime of affinity-dependent growth rates.


\subsection*{T-cell help provides tunable B-cell response specificity}
The specificity of a response follows from the general relation $\DKL = \log L_0 - S_c$ (Materials and Methods), which compares the clonal entropy of the response with that of the naive repertoire.
Primary responses are diversity-dominated ($\zeta < 1$): their clonal entropy is set by the number of activated lineages, $S_c \approx \log L_\act$, so that $\phi \approx L_\act/L_0$ is the fraction of naive clones within reach of the activation front.
Recall responses are size-dominated ($\hat\zeta > 1$): cells are concentrated in the largest clones, and the clonal entropy is set mainly by $\hat\zeta$; because $\hat\zeta$ is close to 1, it retains a weak dependence on the number of re-engaged lineages (Materials and Methods).
In both cases, $\log L_0$ is the maximum specificity, reached by a monoclonal response of the best clone.

The T-cell help threshold $\pi_c$ sets the reach of the activation front.
At the antigen peak, the activation condition gives $\DG_{\max} = \log[k_\on N_A(t_p)/(\lam_A \pi_c)]$, so the reach decreases logarithmically with $\pi_c$ (see Materials and Methods).
A lower $\pi_c$ lets lower-affinity clones join the response (larger $\DG_{\max}$; red in Fig.~\ref{fig:potency}A), whereas a higher $\pi_c$ excludes them (smaller $\DG_{\max}$; cyan in Fig.~\ref{fig:potency}A).
A stringent threshold therefore increases the specificity of the response by focusing it on high-affinity clones.
Because $\pi_c$ drops out of the EP exponent $\alpha$, the clone-size profiles for different $\pi_c$ have the same slope (Fig.~\ref{fig:potency}A): the gain in specificity comes from fewer low-affinity clones taking part in the response, not from larger high-affinity clones.

The coupling of division to T-cell help has a second, independent effect: it sharpens the EP exponent itself. 
To see this, consider the limit of fast pMHC surface turnover, $\delta_\pi \gg \lambda_B$. 
Then $\pi$ tracks the external antigen at all times, $\pi \propto N_A(t)\,\psi(E)$, and dilution by cell division no longer affects termination: clonal expansion stops only as the antigen clears, with termination velocity $v_\text{off} = \delta_A$. 
This limit corresponds to a null model without T-cell help, in which B-cell proliferation is driven directly by antigen capture (see SI), and gives the exponent $\alpha^o = \lambda_B (1/\lambda_A + 1/\delta_A)$. In the dilution-limited regime of the full model ($\lambda_B < \delta_A$), $\alpha^o < \alpha$: without $\pi$-mediated help, the clone-size hierarchy is less steep. 
By coupling division to an internal variable $\pi(t)$ rather than to the external antigen concentration $N_A(t)$, T-cell help converts a clearance-limited termination into a dilution-limited one. T-cell help thus shapes specificity in two ways: the threshold $\pi_c$ sets which clones take part in the response, and the coupling through $\pi$ sets how strongly clone size is graded by affinity.

Additionally, $\pi_c$ sets the duration of clonal expansion.
A higher $\pi_c$ delays activation and accelerates termination for every clone, shortening the proliferation window $\Delta t(\DG) = (\DG_{\max}-\DG)(1/v_\on + 1/v_\off)$ (Eq.~\ref{eq:Delta_t_E}).
As a result, the total primary yield decreases as a power of the threshold, $Y(\pi_c) \sim \exp[\beta^*\DG_{\max}] \sim \pi_c^{-\beta^*}$ (see Materials and Methods).
With $\beta^* \approx 2.3$, a tenfold increase in $\pi_c$ reduces the yield about 200-fold.
A fraction $\nu\gamma$ of this yield is kept as memory.
Although the B-cell energy distribution after a primary response, $q_1(\DG)$, is more uniform than that of the naive repertoire, $q_0(\DG)$ (colored and grey in Fig.~\ref{fig:potency}A, respectively), it is still dominated by the many low-affinity clones near the edge $\DG_{\max}$, and the best clones make up only a small fraction $N_1/Y$ of the primary yield (see Materials and Methods).
In the saturated regime ($\nu \ll 1$, $\gamma Y \gg M$), even the best clones are lost from memory if $M N_1/Y < 1$ (red in Fig.~\ref{fig:potency}A).
In the yield-limited regime ($\nu \approx 1$, $\gamma Y \ll M$), the memory pool is no longer capped by $M$ and tracks the primary yield (cyan in Fig.~\ref{fig:potency}A).

The effect on the memory pool is governed by the interplay of specificity and yield, both controlled by $\pi_c$.
The potency of the recall response is therefore non-monotonic in $\pi_c$ (Fig.~\ref{fig:potency}C,D).
In the saturated regime, increasing $\pi_c$ reduces $\DG_{\max}$ and removes the low-affinity tail of the primary response: the memory pool stays full, but the share of the best clones, $N_1/Y$, grows, which increases potency.
Once $\gamma Y$ falls below $M$, the pool can no longer be filled, and every further increase in $\pi_c$ removes memory cells; the loss of yield then outweighs the gain in affinity.
The optimum $\pi_c^*$ lies at the crossover between these regimes, where the primary yield just saturates the memory pool, $\hat Y_o(\pi_c^*) \sim M$ (see Materials and Methods).
Below $\pi_c^*$, potency is specificity-limited; above it, yield-limited.
}


%::::::::::::::::::::::::::::::::::::::::::::::::::::::::::::::::::::::::::::
\section*{Discussion}
%::::::::::::::::::::::::::::::::::::::::::::::::::::::::::::::::::::::::::::


EP identifies a layer of immune regulation that acts before affinity maturation and operates by a different rule.
In EP, all activated clones divide at the same rate; what binding affinity controls is the \emph{duration} of clonal expansion, bracketed by two fronts in binding energy space whose velocities are set by the growth and the clearance of the antigen.
The selection of high-affinity clones is thus kinetic and it does not require differential fitness between clones.
After an immune response, EP establishes a hierarchy of clone sizes, $N_B \sim \exp{[-\alpha \DG]}$, where the exponent $\alpha$ is set by three macroscopic rates ($\lam_A, \lam_B$ and $\delta_A$) and is independent of the T-cell help threshold ($\pi_c$), the reach of the response ($\DG_{\max}$), and the details of the energy model (e.g, $\beta^*$).
This parameter independence underlies the experimentally observed ratio $\hat\zeta/\zeta = 2$ between recall and primary clone-size exponents.

EP matters because this early phase conditions the three downstream outcomes of the response.
First, it sets the input to AM: the clonal composition entering germinal centers is already graded by affinity with exponent $\alpha$, so the diversity available at the start of the GC reaction is not the diversity of the naive repertoire but a strongly reshaped sample of it.
Second, it sets the properties of the memory compartment.
Because a large fraction of memory is generated with little or no AM~\cite{tangye_human_2007, dogan_multiple_2009, Viant_2020_Fate, Mesin_2020_Restricted_clonality, Glaros_2021_Early_memory}, memory inherits the EP grading of the primary expansion.
Combined with the architecture of the naive repertoire, we predict that the memory compartment is dominated by many clones of suboptimal affinity.
Third, it sets the potency of recall.
The same readout acting a second time on an already graded pool makes recall responses size-dominated rather than diversity-dominated, which concentrates the real-time antibody response on a few high-affinity clones and leaves most memory lineages unengaged, as observed in prime-boost experiments~\cite{Mesin_2020_Restricted_clonality}.
The recall bottleneck and the protective potency of recall are therefore two faces of the same mechanism: the more potent the recall, the earlier the control time $t_c$ and the narrower the set of re-engaged clones.

These three roles are coupled through a single control parameter, the T-cell help threshold $\pi_c$, which identifies and quantifies T-cell help as an important regulator of the B-cell response.
A stringent threshold raises specificity by excluding low-affinity clones, but lowers yield, and hence the number of cells available to fill the memory pool; the potency of recall is maximal at the crossover, where the primary yield just saturates the pool.
Since $\pi_c = b_B/h_0$, the host can shift this optimum by modulating $h_0$, the effective rate of acquisition of T cell help, through the size of the cognate T cell compartment, the co-stimulatory requirements, or the density of pMHC needed to engage a T cell.
This makes the specificity-yield trade-off a property of the immune response subjected to evolutionary tuning.

EP is a non-equilibrium discrimination mechanism, and the analogy with molecular kinetic proofreading~\cite{hopfieldKineticProofreading1974, ninioKineticAmplification1975, mckeithanKineticProofreadingTcell1995} is instructive but incomplete.
Classical proofreading amplifies affinity differences through a cascade of irreversible molecular steps, each paid for in free energy, and acts at the level of a single receptor-ligand complex.
EP amplifies the same differences at the clone level, using the finite duration of antigen availability, bracketed by growth and clearance, as the discriminating variable.
The amplification factor is not a number of molecular steps but a ratio of rates.
Coupling division to the internal variable $\pi(t)$ rather than to the external antigen load is what converts a clearance-limited termination into a dilution-limited one, and sharpens the exponent beyond what antigen dynamics alone would produce.

The framework presented here has a broad impact beyond single immune dynamics. 
It open a new observing windows into the architecture of the immune repertoire.
The clone-size statistics give access to $\beta^*$, the exponent of the naive repertoire density at its high-affinity edge.
Unlike the rates entering $\alpha$, $\beta^*$ is not a property of a given infection but of the repertoire itself.
It summarizes how the number of lineages grows as binding is relaxed.
An evolutionary parameter of the repertoire is thus inferable from a clone-size measurement.

Moreover, we introduce a new way to estimate the specificity of the immune response.
We defined it at the phenotypic level, as the relative entropy between the responding and the naive binding-energy distributions.
This definition is independent of how a clone is delimited, and it quantifies the effective fraction of the naive repertoire that a given antigen engages.
It therefore provides a measurable route to the long-standing question of how many independent responses a repertoire can mount, posed first in the shape-space picture for antigen-antibody pairs~\cite{perelsonOsterTheoreticalStudies1979, Perelson_1997_Review}.


Finally, beyond the immune system itself, our model provide a direct connection with epidemiology.
Epidemiological models encode immune memory through phenomenological cross-immunity functions, which are rarely derived from immunological mechanism~\cite{gogGrenfellDynamicsMultistrain2002, grenfellUnifyingEpidemiologicalEvolutionary2004}.
The potency $P$ defined here is a mechanistic candidate for such a function.
It is a log-fold reduction in peak viral load, and it follows from the structure of the memory repertoire and from the reach of the recall activation front.
Extending the model to a challenge antigen at antigenic distance $d$ from the priming antigen, the calculated potency would give a direct measurement of the breadth and depth (zero distance) of the memory response.
A cross-immunity function of this form could be used in multi-strain epidemiological and predictive models~\cite{fonvilleAntibodyLandscapes2014, luksza_predictive_2014, yanShraimanNeher2019, marchiAntigenicWaves2021}.


\section*{Materials and Methods}

\subsection*{Mean-field equations}
The dynamics is governed by
\EQA
\dot N_A(t) &=& \left[b_A(t) - \delta_A\right] N_A(t), \label{eq:NA}\\
\dot \pi(t, \DG) &=& u_p(t, \DG) - \delta_\pi \,\pi(t, \DG)\nonumber\\
&& - \lam_B(t, \DG)\,\pi(t, \DG)  \label{eq:pi}\\
\dot N_B(t, \DG) &=& \lam_B(t, \DG)\, N_B(t, \DG). \label{eq:NB}
\EEA
Eq.~\ref{eq:NA} describes exponential antigen growth during acute infection; at early times, the antigen sensed by B cells after vaccination follows the same effective dynamics~\cite{MoranTovar_2024_Primary}.
From the onset ($t=0$) until the antigen peak ($t=t_p$), the antigen replicates at a constant birth rate $b_A$ and grows at rate $\lam_A \equiv b_A - \delta_A$; after the peak ($t>t_p$), replication stops and the antigen decays at rate $\delta_A$.
Eq.~\ref{eq:pi} describes the balance of pMHC on the B-cell surface: gain through capture and processing at rate $u_p$, and loss through two channels: surface turnover at rate $\delta_\pi$ and dilution by cell division at rate $\lam_B$, which redistributes the existing pMHC pool into two daughter cells at each division.
Eq.~\ref{eq:NB} describes clonal expansion of B cells at the effective division rate $\lam_B$ given by Eq.~\ref{eq:lambdaB}.

\subsection*{Minimal model of viral fitness}
Assuming binding equilibrium between free antigen and secreted antibody, the probability that a viral particle is free is $(1+Z)^{-1}$, where $Z$ is the antibody-level partition function defined in the Model section (the calibration constant $c_o$ is derived in the SI).

Taking neutralization to block productive infection,
Eq.~\ref{eq:NA} becomes
\EQ
\dot N_A(t) = \left[\frac{b_A}{1 + Z(t)} - \delta_A\right] N_A(t) .
\label{eq:NA_coupled}
\EE
For $Z = 0$ this recovers $\lam_A = b_A - \delta_A$. Net growth vanishes at
the \emph{critical partition function}
\EQ
Z_c = \frac{b_A}{\delta_A} - 1 = \frac{\lam_A}{\delta_A} ,
\label{eq:Zc}
\EE
equivalently $Z_c = R_0 - 1$. Thus $Z_c$ is a property of the pathogen alone,
fixed by the same two rates that set the antigen dynamics in
Fig.~\ref{fig:ep}A; it is not a free parameter.

We note that $c_o$ and $Z_c$ enter the first-passage condition only through the
ratio $Z_c/c_o$ and are therefore not separately identifiable from response data.

\subsection*{Activation and termination fronts}
Before activation, $\lam_B \approx 0$, and for exponential antigen growth, $N_A(t) = N_A(0)\,e^{\lam_A t}$, the pMHC density tracks the source,
\EQ
\pi(t, \DG) \approx \frac{k_\on\, N_A(t)\, \psi(\DG)}{\lam_A},
\label{eq:pi_E}
\EE
where we neglected slow surface turnover, $\delta_\pi \ll \lam_A$ (SI, Moving front).
The activation condition $\pi(t, \DG_\on(t)) = \pi_c$ is therefore equivalent to $\psi(\DG_\on(t)) = \pi_c\, \lam_A / (k_\on\, N_A(t))$.
In the regime $\psi(\DG) \approx e^{-\DG}$, this gives a linear front, $\DG_\on(t) = \lam_A (t - t^*)$, with velocity $v_\on \equiv \partial_t \DG_\on = \lam_A$; the explicit form of $t^*$ is given in the SI.
The cumulative number of naive lineages recruited by time $t \leq t_p$ follows from integrating the lineage density over the swept region of energy space,
\EQ
L_\act(t) \approx \frac{\omega_0}{\beta^*} \exp[\beta^* \DG_\on(t)].
\EE
At the antigen peak, the front stalls at $\DG_{\max} = \log[k_\on N_A(t_p)/(\lam_A \pi_c)]$, which decreases logarithmically with $\pi_c$.

After activation, dilution by division reduces the pMHC density of activated clones to $\pi \approx k_\on N_A(t)\,\psi(\DG)/(\lam_A + \lam_B)$, so that at the peak $\pi_p(\DG) \approx \frac{\lam_A}{\lam_A+\lam_B}\,\pi_c\, e^{\DG_{\max} - \DG}$.
After the peak, Eq.~\ref{eq:pi} has two decay channels: dilution by cell division at rate $\lam_B$, and decay of the antigen source at rate $\delta_A$.
The long-time dynamics is set by the slower of the two (SI, Termination),
\EQ
\pi(t, \DG) \approx K\,\pi_p(\DG)\, e^{-v_\off (t - t_p)}, \qquad v_\off = \min(\delta_A, \lam_B),
\EE
with $K = [\lam_A + \max(\delta_A, \lam_B)]/|\lam_B - \delta_A|$.
When $\delta_A > \lam_B$, the source vanishes quickly and dilution by division dominates the decay; when $\delta_A < \lam_B$, $\pi$ tracks the slowly decaying antigen.
Solving the termination condition $\pi(t_\off, \DG) = \pi_c$ gives Eq.~\ref{eq:t_off}, up to a $\DG$-independent shift $c_0/v_\off$ with $c_0 = \log[K\lam_A/(\lam_A + \lam_B)]$.

Combining Eqs.~\ref{eq:t_on} and~\ref{eq:t_off}, and using $\DG_{\max} = v_\on (t_p - t^*)$, the proliferation window of a clone with binding energy $\DG$ is
\EQ
\Delta t(\DG) = (\DG_{\max} - \DG)\left(\frac{1}{v_\on} + \frac{1}{v_\off}\right)
\label{eq:Delta_t_E}
\EE
up to the same constant, which is linear in $\DG$ and gives Eq.~\ref{eq:alpha}; the constant multiplies all clone sizes by a common factor and does not affect $\alpha$.

\subsection*{Memory formation and sampling}
Each B cell at the end of the primary response becomes memory with probability $p = \nu\gamma$, where $\nu = M/(\gamma Y + M)$ accounts for the carrying capacity $M$ of the memory pool.
The number of memory cells from a clone of size $N_B$ is then binomial with mean $p N_B$, and the expected energy density of memory cells is a scaled copy of the primary one.

Although the rescaling is uniform at the level of cells, its effect on the clonal structure is not: because low-affinity clones are smaller, they are more likely to be lost.
A clone is represented in memory with its mean-field size if $p N_B(\DG) \gtrsim 1$.
With Eq.~\ref{eq:alpha} and $N_B^* \sim \exp[\alpha\DG_{\max}]$, this gives a low-affinity cutoff
\EQ
\DG \lesssim \DG_{\rm mem} \equiv \DG_{\max} - \frac{1}{\alpha}\log\frac{1}{p} = \frac{1}{\alpha}\log\left(p N_B^*\right).
\label{eq:E_mem}
\EE
Clones beyond the cutoff survive with probability $\approx p N_B(\DG)$, typically as single cells.
Memory formation thus preserves the exponential profile of the high-affinity clones and flattens the low-affinity tail.

The total primary yield follows from integrating the clone sizes over the activated repertoire,
\EQA
Y &=& \int_{\DG^*}^{\DG_{\max}} \Omega_0(\DG)\, N_B(\DG)\, d\DG\nonumber\\
&\approx& \frac{\omega_0}{\beta^*(1-\zeta)}\, \exp[\beta^*\DG_{\max}] = \frac{L_\act}{1-\zeta},
\label{eq:yield}
\EEA
where we used $\DG^* \approx 0$ and $\beta^*(1-\zeta)\DG_{\max} \gg 1$, valid for $\zeta < 1$.
The yield is thus dominated by the many low-affinity clones near $\DG_{\max}$, and the best clone makes up only a fraction
\EQ
\frac{N_B^*}{Y} \approx (1-\zeta)\, \frac{N_B^*}{L_\act} \sim \exp[-\beta^*(1-\zeta)\DG_{\max}]
\EE
of it.

The best clone is retained if $\DG_{\rm mem} \gtrsim \DG^*$, that is, if $p N_B^* \gtrsim 1$.
In the saturated regime ($\gamma Y \gg M$, $p \approx M/Y$), this requires $M N_B^*/Y \gtrsim 1$, which fails if the yield is spread over many low-affinity clones.
In the unsaturated regime ($\gamma Y \ll M$, $p \approx \gamma$), it requires $\gamma N_B^* \gtrsim 1$, which fails if the best clone is small.
If the condition fails, sampling also thins the high-affinity edge.
The best affinity represented in memory is then set by the $pY$ sampled cells: with the primary cell density $\propto \exp[\beta^*(1-\zeta)\DG]$, it lies at
\EQ
\hat\DG^* \approx \DG_{\max} - \frac{\log(pY)}{\beta^*(1-\zeta)}.
\EE

\subsection*{Recall response}
Upon re-challenge, memory clones are activated and terminated by the same mechanism as in the primary response, with unchanged $\psi(\DG)$, $k_\on$ and $\lam_B$, but with a lower activation threshold $\hat\pi_c < \pi_c$.
The two fronts therefore advance with the same velocities, and the proliferation window of a memory clone is $\Delta\hat t(\DG) = \hat C - \DG\,(1/v_\on + 1/v_\off)$, where $\hat C$ is independent of $\DG$ (SI, Memory response).
With the inherited abundance $\hat N_B^o(\DG) = \nu\gamma N_B^*\, e^{-\alpha\DG}$, the final size of a memory clone is $\hat N_B(\DG) = \hat N_B^o(\DG)\, e^{\lam_B \Delta\hat t(\DG)} \propto e^{-2\alpha\DG}$, which gives Eq.~\ref{eq:alpha2}.

We assume that the recall challenge starts from the same initial antigen load and grows at the same rate $\lam_A$ as the primary one.
The recall front then stalls at the control time $t_c$, where $N_A(t_c)/N_A(t_p^o) = e^{-P}$ (Eq.~\ref{eq:P}).
Applying the activation condition with threshold $\hat\pi_c$ at $t_c$ and with $\pi_c$ at $t_p^o$, the kinetic prefactors cancel, and the reach of the recall front is
\EQ
\hat\DG_{\max} = \DG_{\max} - P + \Delta,
\qquad
\Delta \equiv \log\frac{\pi_c}{\hat\pi_c}.
\label{eq:Emax_recall}
\EE
Since memory lineages are distributed with density $\propto e^{\beta^*\DG}$ up to $\DG_{\max}$, the fraction of primary lineages that are re-engaged is
\EQ
\frac{\hat L_\act}{L_\act} \approx e^{-\beta^*(P - \Delta)},
\label{eq:bottleneck}
\EE
valid for $\hat\DG_{\max} < \DG_{\rm mem}$ (Eq.~\ref{eq:E_mem}).
The re-engaged fraction of memory cells, the self-limitation of the bottleneck, and its cost in potency are derived in the SI (Recall bottleneck).

\subsection*{Clone-size statistics}
With $N_B(\DG) \propto e^{-\alpha\DG}$ (Eq.~\ref{eq:alpha}) and $\Omega_0(\DG) = \omega_0 e^{\beta^*\DG}$, the rank of a clone with binding energy $\DG$ is $k(\DG) \approx (\omega_0/\beta^*)\,e^{\beta^*\DG}$, so that clone sizes follow a Zipf law, $N_k \propto k^{-\zeta}$ with $\zeta = \alpha/\beta^*$; for recall responses, $\alpha \to 2\alpha$ gives $\hat\zeta = 2\zeta$ (Eq.~\ref{eq:zetas}).
The largest clone and the number of activated lineages are set by the reach of the activation front, $N_1 \sim e^{\alpha\DG_{\max}}$ and $L_\act \approx (\omega_0/\beta^*)\,e^{\beta^*\DG_{\max}}$, which gives $L_\act \sim N_1^{1/\zeta}$ (Eq.~\ref{eq:scaling}).
The total yield, $Y = \omega_0 e^{\alpha\DG_{\max}} \int_0^{\DG_{\max}} e^{(\beta^* - \alpha)\DG}\,\d\DG$, is dominated by clones near $\DG_{\max}$ for $\zeta < 1$ (Eq.~\ref{eq:yield}) and by clones near $\DG^*$ for $\zeta > 1$, where $Y \approx \omega_0 e^{\alpha\DG_{\max}}/(\alpha - \beta^*)$.
If each cell is sequenced with probability $p_s$, clones with $p_s N_k \gtrsim 1$ keep their relative sizes, so the shape of normalized rank curves is unchanged, and the scaling relation keeps its exponent (SI, Clone-size scaling laws).

In the rate model, $\lam_B \approx h_0\pi$ gives Eq.~\ref{eq:rate_model}, and hence $\log N_k = \log N_1\, k^{-c}$ with $c = 1/\beta^*$.
In normalized form, $\log(N_k/N_1) = z_0\,(k^{-c} - 1)/c$ with $z_0 = c\log N_1$, which reduces to the Zipf law with $\zeta = z_0$ for $c \to 0$; $N_1$ is a fit parameter because the sampling fraction is unknown (SI, Affinity-dependent expansion rates).

\subsection*{Repertoire exponent in the Gaussian approximation}
In the Gaussian approximation, the naive energy density has mean $\mu_0$ and variance $\sigma^2$, truncated at the best clone $\DG^*$.
By extreme-value statistics, the best of $L_0$ clones lies at $\Delta\DG \equiv \mu_0 - \DG^* \approx \sigma\sqrt{2\log L_0}$, and the local exponent of the density at $\DG^*$ is $\beta^* = \Delta\DG/\sigma^2$.
Hence $\beta^* = \sqrt{2\log L_0}/\sigma$.
For a binding energy that is a sum of $\ell$ independent residue contacts, $\sigma^2 = \ell\sigma_r^2$, and with $\sigma_r = 1$ this gives $\beta^* = \sqrt{2\log L_0/\ell}$.


\subsection*{Specificity and clonal entropy}
The fraction of B cells belonging to clone $i$ is $p_i$ such that $\sum_i^{L_0}p_i = 1$.
We define the clonal entropy as the Shannon entropy
\EQ
S_c\left[\{p_i\}\right] = -\sum_i^{L_0} p_i \log p_i
\EE
For the naive repertoire, then $p^0_i=1/L_0$ and the clonal entropy of the naive repertoire is given by 
\EQ
S^0_c\equiv S_c\left[\{p^0_i\}\right]=\log L_0.
\EE

At the phenotypic level, the continuous entropy of the distribution $q$ is given by
\EQ
S_\DG\left[q\right] = -\int q(\DG)\log q(\DG) \, \d \DG 
\EE
For the naive repertoire, the corresponding energy pdf is given by $q_0(\DG) = \Omega_0(\DG)/L_0$.
In this case, the entropy is given by
\EQA
S_E^0\equiv S_\DG\left[q_0\right] &=& \log L_0 -\left\langle \log\Omega_0(E)\right\rangle_{q_0}\nonumber\\
&=& S_c^0 -\left\langle \log\Omega_0(E)\right\rangle_{q_0}
\EEA

In general, for any arbitrary clone-size distribution $p_i$ and energy distribution $q(\DG)$, the following relation holds
\EQ
S_c\left[\{p_i\}\right] = \left\langle \log\Omega_0(\DG)\right\rangle_{q} + S_\DG\left[q\right]
\label{eq:S_c_S_E}
\EE
The proof, which assumes that clone size is a deterministic function of binding energy, is given in the SI (Specificity and diversity).

We also calculate the relative entropy between the distribution $q(\DG)$ and the naive distribution $q_0(\DG)$
\EQA
\DKL\left[q||q_0\right] &=& \int q(\DG)\log \frac{q(\DG)}{q_0(\DG)} \, \d \DG\nonumber\\
&=& -S_\DG\left[q\right] - \left\langle \log\Omega_0(\DG)\right\rangle_{q} + \log L_0\nonumber\
\EEA
Using \ref{eq:S_c_S_E}, we obtain
\EQ
\DKL\left[q||q_0\right] = \log L_0 - S_c\left[\{p_i\}\right]
\EE
Interestingly, this last equation is also the Kulback-Leibler distance between the clone-size distribution $p_i$ and the uniform distribution $p_i^0$
\EQA
\DKL\left[\{p_i\}||\{p_i^0\}\right] &=& \sum_i^{L_0} p_i\log \frac{p_i}{p_i^0}\nonumber\\
&=& -S_c\left[\{p_i\}\right] + \log L_0
\EEA
We calculate the fraction of the activated repertoire, $\phi$, as
\EQA
\phi &\equiv& \exp\left(-\DKL\right)\nonumber\\ 
&=& \frac{1}{L_0}\exp\left(S_c\left[\{p_i\}\right]\right)
\EEA

For Zipf-distributed clone sizes, $p_k = k^{-\zeta}/H$ with $H = \sum_{k=1}^{L_\act} k^{-\zeta}$, the continuum approximation gives
\EQ
S_c(\zeta, L_\act) = \log\frac{1-x}{\zeta-1} + \frac{\zeta}{\zeta-1}\,\frac{1-x(1+y)}{1-x},
\label{eq:Sc_zipf}
\EE
with $x \equiv L_\act^{1-\zeta}$ and $y \equiv (\zeta-1)\log L_\act$, valid for either sign of $\zeta - 1$.
For a diversity-dominated response ($\zeta < 1$, $(1-\zeta)\log L_\act \gg 1$), $S_c \approx \log L_\act - \zeta/(1-\zeta) - \log(1-\zeta)$.
For a size-dominated response ($\zeta > 1$), $S_c \to \zeta/(\zeta-1) - \log(\zeta-1)$ as $L_\act \to \infty$; for the measured $\hat\zeta \approx 1.13$ and $L_\act$ of order $50$, $x \approx 0.6$ and the finite form, Eq.~\ref{eq:Sc_zipf}, must be used (SI).


\subsection*{Dependence on the T-cell help threshold}

At the antigen peak, the activation front stalls at
\EQ
\DG_{\max} = \log\frac{\pi_p^*}{\pi_c}, \qquad \pi_p^* \equiv \frac{k_\on N_A(t_p)}{\lam_A},
\label{eq:Emax_pic}
\EE
where $\pi_p^*$ is the peak pMHC density of the best clone.
In primary responses, the antigen peak is set independently of the B-cell response (by innate immunity in an infection), so $N_A(t_p)$, and hence $\pi_p^*$, does not depend on $\pi_c$.
The threshold therefore enters the response only through $\DG_{\max}$.
With Eq.~\ref{eq:yield} and Eq.~\ref{eq:alpha}, the primary yield and the size of the best clone are
\EQ
Y \approx \frac{\omega_0}{\beta^*(1-\zeta)} \left(\frac{\pi_p^*}{\pi_c}\right)^{\beta^*}, \qquad
N_B^* \approx \left(\frac{\pi_p^*}{\pi_c}\right)^{\alpha},
\EE
valid as long as $\beta^*(1-\zeta)\DG_{\max} \gg 1$.
Increasing $\pi_c$ thus reduces both, but the yield faster, so the share of the best clone, $N_B^*/Y \sim (\pi_c/\pi_p^*)^{\beta^*(1-\zeta)}$, grows.

The potency of the recall response is maximal where the primary yield slightly overfills the memory pool,
\EQ
\gamma Y(\pi_c^*) = \frac{s^*}{1-s^*}\, M, \qquad s^* = \hat\zeta - \frac{1}{\beta^*},
\EE
which balances the delayed activation and smaller best clone at higher thresholds against the purging of the low-affinity tail (SI, The T-cell help threshold and the optimal potency).
With the measured exponents, $s^* \approx 0.65$--$0.85$, so the optimum lies close to the crossover from the saturated to the unsaturated regime.
Using the expression for $Y$ above,
\EQ
\pi_c^* \approx \pi_p^* \left[\frac{(1-s^*)\,\gamma\,\omega_0}{s^*\, M \beta^*(1-\zeta)}\right]^{1/\beta^*}.
\EE

The best clone is retained in memory if $\nu\gamma N_B^* \gtrsim 1$ (see Memory formation and sampling).
In the saturated regime, this requires $M N_B^*/Y \gtrsim 1$, which holds for sufficiently large $\pi_c$; in the unsaturated regime, it requires $\gamma N_B^* \gtrsim 1$, which holds for $\pi_c \lesssim \pi_p^*\gamma^{1/\alpha}$.
The best clones are therefore retained within a window of thresholds: at low $\pi_c$, the memory pool fills with low-affinity cells, and at high $\pi_c$, the best clone is too small to survive sampling.

\subsection*{Repertoire data analysis}
\todo{Datasets and sorts used (CGG-alum experiments pooled; influenza/HA), identity of the HA ``GC'' sort, GC vs.\ plasmablast compartments, clone definition (V, D, J), normalization by the largest sampled clone, fixed random seed.}

We fitted the Zipf law, Eq.~\ref{eq:SI_rate_zipf_norm}, and the two-parameter family, Eq.~\ref{eq:SI_rate_family} with $c \geq 0$, to the cohort-averaged normalized rank curves of Fig.~\ref{fig:csd}B.
At each rank, $x_k$ was averaged over the mice with at least $k$ clones; ranks with fewer than two contributing mice were discarded, and fits used ranks $2 \leq k \leq 20$ (rank 1 is fixed by normalization).
Both forms were fitted by least squares on $\log x_k$, and compared by
\EQ
\Delta{\rm AIC} = {\rm AIC}_{\rm rate} - {\rm AIC}_{\rm Zipf},
\qquad
{\rm AIC} = n \log\frac{\rm RSS}{n} + 2k,
\label{eq:rate_AIC}
\EE
where RSS is the residual sum of squares, $n$ the number of fitted ranks, and $k$ the number of parameters ($k = 1$ for Zipf, $k = 2$ for the rate model).
If the curvature parameter does not improve the fit, the rate model converges to $c = 0$ and $\Delta{\rm AIC} = +2$.
Uncertainties were obtained by resampling mice with replacement ($400$ bootstrap replicates); since residuals at consecutive ranks are correlated, we use the bootstrap distribution of $c$ rather than absolute AIC values as the primary measure.

\subsection*{Numerical simulations}
\todo{Numerical scheme (one paragraph); parameter table in SI (Numerical simulations).}


\section*{Acknowledgments}
This work was partially supported by the Deutsche Forschungsgemeinschaft (grant CRC 1310, to M.L.).

\bibliographystyle{apsrev4-2}
\bibliography{references}
  

\todo{
\begin{itemize}
	\item In the Gaussian approximation, the following relations hold:
	\begin{enumerate}
		\item $\Delta \DG \equiv \mu_0 - \DG^*$
		\item $\Delta \DG = \beta^*\sigma^2$
		\item $\Delta \DG = \sigma\sqrt{2\log L_0}$
		\item $\mu_0 =  \ell \mu_r$ ; $\sigma^2 = \ell \sigma_r^2$
		\item values: $\beta^*\approx 2.3$ from data.
		\item if $\mu_r \approx 0$, then $\mu_0=0$
		\item for $L_0 = $
		\item $K_D^* \approx 0.1\,\mu M$ (from UCA from \cite{Mesin_2020_Restricted_clonality})
		\item Average affinity after a primary response $K_D \approx 10 \,\mu M$ (from \cite{Eisen_2014_AM, AltanBonnet_2020_Review}) 
	\end{enumerate}
	where $\mu_0$ is the energy distance between the average energy and the EVT value. 
\end{itemize}
}

\end{document}

